\chapter{Auditoría: verificar y anclar}
\label{cap:auditoria}

El registro de auditoría es una cadena de hashes: cualquier cambio en un evento rompe la cadena.
Pero alguien con acceso de administrador a MySQL podría borrar los eventos más recientes o
reescribir la cadena completa. Para detectarlo, se \textbf{ancla} periódicamente la cabeza de la
cadena en un archivo guardado fuera de su alcance: el USB del auditor o una carpeta compartida
de otra PC en la que el administrador de la base no pueda escribir.

\begin{xltabular}{\linewidth}{@{}l L@{}}
	\toprule
	\textbf{Cuándo} & \textbf{Qué} \\
	\midrule
	\endhead
	Diario & Verificar la cadena y las anclas \\
	Semanal y al cerrar cada proceso de admisión & Anclar la cabeza actual \\
	Tras cada restauración, actualización o uso de la cuenta \texttt{root} de MySQL & Verificar
	y guardar la salida con el registro de la intervención \\
	\bottomrule
\end{xltabular}

\section{Comandos}

\so{Docker} (desde \texttt{deploy}; el archivo de anclas está en \texttt{ANCHOR\_DIR})
\begin{codigo}
docker compose --profile ops run --rm audit
docker compose --profile ops run --rm audit anchor --anchors /anchors/digaipe-anchors.jsonl --note "cierre UA06426"
\end{codigo}

\so{Linux} (con el USB del auditor montado, por ejemplo, en \texttt{/media/auditor})
\begin{codigo}
sudo /usr/local/sbin/digaipe-audit verify --anchors /media/auditor/digaipe-anchors.jsonl
sudo /usr/local/sbin/digaipe-audit anchor --anchors /media/auditor/digaipe-anchors.jsonl --note "cierre UA06426"
\end{codigo}

\so{Windows} (con el USB del auditor, por ejemplo, en \texttt{E:})
\begin{codigo}
powershell -ExecutionPolicy Bypass -File C:\digaipe\bin\digaipe-audit.ps1 verify --anchors E:\digaipe-anchors.jsonl
powershell -ExecutionPolicy Bypass -File C:\digaipe\bin\digaipe-audit.ps1 anchor --anchors E:\digaipe-anchors.jsonl --note "cierre UA06426"
\end{codigo}

\texttt{anchor} verifica primero la cadena y se niega a anclar una cadena rota. Imprime la línea
del ancla: cópiela también en un registro independiente (libro de actas, correo interno). Cada
ancla queda además en la propia cadena como evento \texttt{audit.anchored}.

\begin{aviso}[Si la verificación falla]
	\textbf{No repare nada.} Conserve la base, los respaldos y el archivo de anclas tal como están
	y registre el incidente. La salida indica qué eventos no coinciden. Si dice
	\texttt{audit triggers missing}, ejecute la instalación (\texttt{setup}) y verifique de nuevo;
	no deje la aplicación funcionando sin triggers.
\end{aviso}

Después de una restauración, la verificación con \texttt{--anchors} marcará como faltantes las
anclas posteriores al respaldo restaurado. Es lo esperado: el evento \texttt{backup.restored} y el
respaldo de seguridad lo documentan. Guarde esa salida con el registro del incidente.

\section{Cuenta de solo lectura para auditores}

Si los auditores necesitan consultar la base con herramientas SQL, defina
\texttt{READONLY\_DB\_PASSWORD} (en \texttt{deploy/.env} o en \texttt{setup.env}) y vuelva a
ejecutar la instalación. La cuenta \texttt{digaipe\_lectura} lee todas las tablas, incluida la
auditoría, pero no ve contraseñas, tokens de sesión ni códigos de verificación, no puede
modificar nada y exige conexión cifrada (TLS). Como MySQL no acepta conexiones desde la red, se
usa desde el propio servidor:

\begin{xltabular}{\linewidth}{@{}l L@{}}
	\toprule
	\textbf{Modo} & \textbf{Comando} \\
	\midrule
	\endhead
	Docker & \texttt{docker compose exec db mysql -h 127.0.0.1 --ssl-mode=REQUIRED -u
	digaipe\_lectura -p digaipe} \\
	Linux / Windows & \texttt{mysql -h 127.0.0.1 --ssl-mode=REQUIRED -u digaipe\_lectura -p
	digaipe} \\
	\bottomrule
\end{xltabular}
